\scanpage{71}
\mspage{77}

La proposition duale fait appel aux foncteurs can
\[
{}_{y_0\backslash}X\longrightarrow {}_{y_1\backslash}X
\qquad\text{et}\qquad X_y\longrightarrow {}_{y\backslash}X,
\]
et aux faisceaux $f_! f^*M$, pour $M$ constant ou \unsure{loc|loc.} constant sur $Y$.

\uline{Définition 1} Si les conditions équiv. de la prop. 1 sont satisfaites, on dit que $f$ est une $\mathbb D$-fibration cohomologique \deleted{\unsure{faible}} \inserted{(ou $\mathbb D$-fibration \uline{gauche})}. Si les conditions duales sont satisfaites, on dit que c'est une $\mathbb D$-fibration homologique (ou $\mathbb D$-fibration \uline{droite}).

\uline{NB} \unsure{Les} $\mathbb D$-fibrations cohomo\uline{lo}giques sont stables par tt chgt de base qui \deleted{\unsure{est}} un iso local. Les $\mathbb D$-fibrations homologiques par \deleted{\unsure{tout}} chgt de base qui est un co-iso local. Mais ce n'\deleted{\unsure{est pas une}} sont pas des notions stables par \uline{tout} changt de base. Mais \unsure{la} notion de $\mathbb D$-fibration coh. est stable par changt de base $\mathbb D$-lisse, et celle de $\mathbb D$-fibration homologique par chgt de base $\mathbb D$-propre.

\margin{? donner exemples !}

Je vais essayer de construire un exemple où le chgt de base coule, en prenant $Y=\Delta_1$. Se donner $X$ sur $\Delta_1$. revient à se donner une sous-catégorie ouverte $X_0$, et sa complémentaire fermée $X_1$
\diagram{\begin{tikzcd}[ampersand replacement=\&,column sep=small]
\& X \& \\
X_1 \arrow[ur,phantom,"{\unsure{\subset|-}}" description,sloped] \arrow[dd,dotted,dash] \& \& X_0 \arrow[ul,phantom,"{\unsure{\supset|-}}" description,sloped] \arrow[dd,dashed,dash]\\
\& \& \\
1 \& \& 0 \arrow[ll]
\end{tikzcd}}
Dire que $X\longrightarrow Y=\Delta_1$ est \deleted{coh.} $\mathbb D$-fibr. coh. signifie simplement
% archive sheet number: 70
